\documentclass[10pt,a4paper]{amsart}
\usepackage[margin=19mm]{geometry}
\usepackage{amsmath,amssymb,mathtools}
\usepackage[scr=rsfs,cal=euler]{mathalfa}
\usepackage{xfrac}
\usepackage[unicode,hidelinks,bookmarks=false]{hyperref}
\newcommand{\ie}{i.e.\ }
\newcommand{\cf}{cf.\ }
\newcommand{\resp}{respectively\ }
\newcommand{\Subsection}{\subsection}
\providecommand{\nobreakdash}{\nobreak\hbox{-}}
\setlength{\emergencystretch}{2em}
\multlinegap0pt
\usepackage{bm}
\usepackage{bbm}
\usepackage{dsfont}
\usepackage{xspace}
\usepackage{cancel}
\usepackage{mathtools}
\usepackage[shortlabels]{enumitem}
\usepackage{amsthm}
\makeatletter
\@ifundefined{Th}{\newtheorem{Th}{Th}}{}
\@ifundefined{Lemma}{\newtheorem{Lemma}[Th]{Lemma}}{}
\makeatother
\title[Bordered Legendrian Rational Symplectic Field Theory]{Bordered Legendrian Rational Symplectic Field Theory}
\author{Maciej Wlodek}
\date{Source: arXiv:2509.07921v1 (CC BY 4.0)}
\begin{document}
\maketitle

\noindent\emph{Source and licence.} Bordered Legendrian Rational Symplectic Field Theory. Version arXiv:2509.07921v1 (CC BY 4.0), distributed under the Creative Commons Attribution 4.0 International licence, as recorded in the arXiv licence metadata for this exact version and shown on the abstract page https://arxiv.org/abs/2509.07921v1. Full rights notice: licenses/arxiv-2509.07921-CC-BY-4.0.txt.\par\smallskip

\noindent\textit{Editorial scope.} Counted unit: the Lemma whose source label is \texttt{lem:deltaMderivation} in the article (source0.tex lines 1529--1533, source label \texttt{lem:deltaMderivation}), delivered together with the article's own complete original proof of it (source0.tex lines 1535--1652, one \texttt{proof} environment of 118 lines). No mathematics is written, repaired or re-derived: statement and proof are the article's own text line for line. Required context retained uncounted: none. Every definition, display and label of those lines is retained unchanged. Omitted from this package and not redistributed: 1--1528, 1534--1534, 1653--6403. Nothing inside a retained excerpt is omitted or altered: every retained body is delivered exactly as the article's lines are written, so no omission receipt of any kind accompanies this package. Citations: the retained excerpts carry no citation command at all, so no bibliography entry of the article is transcribed and no reference is redistributed. Reference adaptation: none was needed and none was made; every reference inside the delivered unit resolves inside it, so no unresolved reference or citation is produced. Printed slips kept literal: no printed word, symbol, hypothesis or number of the counted statement or of its proof was altered. Layout and class: the article's own class is \texttt{amsart} with packages including amsmath, amsthm, amsfonts, amssymb, enumitem, hyperref, cleveref, graphicx, tikz-cd, tikz; the wrapper is the lane's pinned \texttt{amsart} editorial wrapper at 10pt on A4, which restates the theorem-like environments the retained text needs and adds the article's own macro definitions that the retained text needs (none), with extra wrapper packages (bm, bbm, dsfont, xspace, cancel, mathtools, enumitem). The delivered numbering is the unit's own. Assets: no retained span contains a figure environment, a graphics-inclusion command, a PDF-page inclusion, a table or a TikZ picture; no image file is copied into this package, the package declares no asset dependency and no image file is redistributed with it. Licence: version arXiv:2509.07921v1 of this article is distributed under the Creative Commons Attribution 4.0 International licence, as recorded in the arXiv licence metadata for this exact version and shown on the abstract page https://arxiv.org/abs/2509.07921v1. A full rights notice is delivered at licenses/arxiv-2509.07921-CC-BY-4.0.txt. Characters: every retained excerpt is pure ASCII, so the delivered engine, XeLaTeX with its default Latin Modern text font, reports no missing character.
	\begin{Lemma}[$\delta_{str}^M$ is a derivation of the SFT bracket] \label{lem:deltaMderivation}
		
		\[ \delta_{str}^M\left(\{x,y\}\right) = \{\delta_{str}^M( x), y\} + \{x, \delta_{str}^M(y)\}.\]
		
	\end{Lemma}

	\begin{proof}
		It suffices to prove it when $x,y$ are generators.
		\begin{itemize}
			\item $x=\beta_{ij}^*$, $y=\beta_{k\ell}^*$.
			
			Both sides vanish so the formula holds trivially.
			
			\item $x=\alpha_{ij}^R$, $y=\beta_{k\ell}^L$. 
			
			Both sides vanish so the formula holds trivially. 
			
						
			\item $x=\alpha_{ij}^R$, $y=\beta_{k\ell}^R$. 
			
			\begin{align*}
				&\{\alpha_{ij}^R, \delta_{str}^M(\beta_{k\ell}^R)\} = \{\alpha_{ij}^R, (\beta_{k\ell}^R)^2\} = 2\{\alpha_{ij}^R, \beta_{k\ell}^R\}\beta_{k\ell}^R = 0
			\end{align*}
			
			So it suffices to show
			
			\[ \delta_{str}^M\left(\{\alpha_{ij}^R,\beta_{k\ell}^R\}\right) = \{\delta_{str}^M(\alpha_{ij}^R), \beta_{k\ell}^R\}\]
			
			Expand the RHS.
			
			\begin{align*}
				\{\delta_{str}^M(\alpha_{ij}^R), \beta_{k\ell}^R\} &= \left\{ (\gamma_i+\gamma_j) \alpha_{ij}^R + \sum_{i<m<j} \alpha_{im}^R\alpha_{mj}^R , \beta_{k\ell}^R \right\} \\
				&= \left\{(\gamma_i+\gamma_j), \beta_{k\ell}^R \right\} \alpha_{ij}^R + (\gamma_i+\gamma_j)\left\{ \alpha_{ij}^R, \beta_{k\ell}^R \right\} + \sum_{i<m<j} \left\{\alpha_{im}^R\alpha_{mj}^R, \beta_{k\ell}^R \right\}
			\end{align*}
			
			The first term vanishes since $\gamma_i$ and $\gamma_j$ are both just a sum of $\beta$ generators. We now split into cases based on $\{k,\ell\} \cap \{i,j\}$.
			
			\begin{enumerate}[label=Case \arabic*.]
				\item $\{k,\ell\} \cap \{i,j\} = \varnothing$.
				
				Then the second term $(\gamma_i+\gamma_j)\left\{ \alpha_{ij}^R, \beta_{k\ell}^R \right\}$ vanishes. 
				
				For the third term, for each $i<m<j$, the only way $\left\{\alpha_{im}^R\alpha_{mj}^R, \beta_{k\ell}^R \right\}$ might be nonzero is if $m\in \{k,\ell\}$. But even then $\left\{\alpha_{im}^R\alpha_{mj}^R, \beta_{k\ell}^R \right\} = 2\alpha_{im}^R\alpha_{mj}^R = 0$. Therefore the third term vanishes and 
				
				\[ \{\delta_{str}^M(\alpha_{ij}^R), \beta_{k\ell}^R\} = 0\]
				
				Likewise, 
				
				\[ \delta_{str}^M\left(\{\alpha_{ij}^R,\beta_{k\ell}^R\}\right) = 0. \]
				
				\item $\{k, \ell\} = \{i,j\}$.
				
				For each $i<m<j$, $\left\{\alpha_{im}^R\alpha_{mj}^R, \beta_{k\ell}^R \right\} = 2\alpha_{im}^R\alpha_{mj}^R = 0$. Everything else also vanishes as in case 1.
				
				\item $|\{k,\ell\} \cap \{i,j\}| = 1$. 
				
				Assume $i=k$ and $\ell\neq j$; the other possibilities are analogous.
				
				Let $i<m<j$. If $\ell \neq m$ then
				
				\begin{align*}
					\left\{\alpha_{im}^R\alpha_{mj}^R, \beta_{k\ell}^R \right\} &= \alpha_{im}^R\left\{\alpha_{mj}^R, \beta_{k\ell}^R \right\} + \left\{\alpha_{im}^R, \beta_{k\ell}^R \right\}\alpha_{mj}^R =  0+\alpha_{im}^R\alpha_{mj}^R = \alpha_{im}^R\alpha_{mj}^R
				\end{align*}
				
				If $\ell = m$ then
				
				\begin{align*}
					\left\{\alpha_{im}^R\alpha_{mj}^R, \beta_{k\ell}^R \right\} &= \alpha_{im}^R\left\{\alpha_{mj}^R, \beta_{k\ell}^R \right\} + \left\{\alpha_{im}^R, \beta_{k\ell}^R \right\}\alpha_{mj}^R =  \alpha_{im}^R\alpha_{mj}^R+0 = \alpha_{im}^R\alpha_{mj}^R
				\end{align*}
				
				Thus we have 
				
				\begin{align*}
					\{\delta_{str}^M(\alpha_{ij}^R), \beta_{k\ell}^R\} &= (\gamma_i+\gamma_j) \alpha_{ij}^R + \sum_{i<m<j} \alpha_{im}^R\alpha_{mj}^R \\
					&= \delta_{str}^M(\alpha_{ij}^R) \\
					&= \delta_{str}^M\left(\{\alpha_{ij}^R,\beta_{k\ell}^R\}\right)
				\end{align*}
				
				as desired.
				
			\end{enumerate}
			
			
			\item $x=\alpha_{ij}^R$, $y=\alpha_{k\ell}^L$. 
			
			\begin{align*}
				\delta_{str}^M\left(\{\alpha_{ij}^R,\alpha_{k\ell}^L\}\right) = 0
			\end{align*}
			
			\begin{align*}
				& \{\delta_{str}^M(\alpha_{ij}^R), \alpha_{k\ell}^L\} + \{\alpha_{ij}^R, \delta_{str}^M(\alpha_{k\ell}^L)\} \\
				&=\left\{(\gamma_i+\gamma_j)\alpha_{ij}^R + \sum_{i<m<j} \alpha_{im}^R\alpha_{mj}^R , \alpha_{k\ell}^L\right\} + \left\{\alpha_{ij}^R, (\gamma_k+\gamma_{\ell}) \alpha_{k\ell}^L + \sum_{k<m<\ell} \alpha_{km}^L \alpha_{m\ell}^L  \right\} \\
				&= \left\{\gamma_i+\gamma_j , \alpha_{k\ell}^L\right\}\alpha_{ij}^R + \left\{\alpha_{ij}^R, \gamma_k+\gamma_{\ell}   \right\}\alpha_{k\ell}^L
			\end{align*}
			
			Introduce an ordering $<'$ of $\{1,\dots, n\}$ according to the sequence 
			\[ 1 <' \beta^R(1) <' \beta^L(\beta^R(1)) <' \beta^R(\beta^L(\beta^R(1))) <' \dots \]
			(stopping just before $1$ is reached again).
			If $ij$ and $k\ell$ are not interlaced with respect to $<'$, then both terms above will vanish. If $ij$ and $k\ell$ are interlaced, e.g. $i<'k<'j<'\ell$, then
			
			\[ \{\gamma_i, \alpha_{k\ell}^L\} = 0; \quad \{\gamma_j, \alpha_{k\ell}^L\} = \alpha_{k\ell}^L \]
			
			and 
			
			\[ \{\alpha_{ij}^R, \gamma_k\} = \alpha_{ij}^R; \quad \{\alpha_{ij}^R, \gamma_\ell\} = 0 \]
			
			
			Thus, 
			
			\[ \left\{\gamma_i+\gamma_j , \alpha_{k\ell}^L\right\}\alpha_{ij}^R + \left\{\alpha_{ij}^R, \gamma_k+\gamma_{\ell}   \right\}\alpha_{k\ell}^L = 0,\]
			
			as desired.
			
			
			
			
			\item All other cases are similar to one of the cases above.
			
			
			
			
			
		\end{itemize}
	\end{proof}

\end{document}
